\documentclass[11pt, a4paper]{article}

\usepackage[utf8]{inputenc}
\usepackage[margin=1in]{geometry} % Balanced margins
\usepackage{parskip}              % Modern spacing between paragraphs (no indent)
\usepackage{xcolor}               % For custom colors
\usepackage{graphicx}             % For images
\usepackage{tikz}
\usepackage{amsmath}
\usepackage{amsfonts}
\usepackage{amssymb}
\usepackage{enumitem}

\definecolor{primary}{RGB}{0, 102, 204} % Professional blue
\usepackage[colorlinks=true, linkcolor=primary, urlcolor=primary, citecolor=primary]{hyperref}

\title{\textbf{\color{primary}STAT5147 - Homework 5}}
\author{Aheer Sravon}
\date{\today}

\begin{document}

\maketitle


\section{Problem 1 Answers}

We are given a population of $N=9$ executives, consisting of
$4$ married, $3$ never-married, and $2$ divorced executives.
A random sample of $n=3$ executives is selected without
replacement.

Let $Y_1$ denote the number of married executives and
$Y_2$ denote the number of never-married executives
in the sample.

Since the selection is made without replacement,
$(Y_1,Y_2)$ follows a multivariate hypergeometric
distribution.

For a hypergeometric random variable $Y$ with population
size $N$, number of successes $K$, and sample size $n$,
the expectation and variance are

\begin{align*}
E(Y) &= n\frac{K}{N},\\
\operatorname{Var}(Y)
&= n\frac{K}{N}
\left(1-\frac{K}{N}\right)
\frac{N-n}{N-1}.
\end{align*}

\subsection{(a) Expectation and Variance of $Y_1$}

There are $K=4$ married executives among the $N=9$
executives, and $n=3$ are selected.

The expected number of married executives is

\begin{align*}
E(Y_1)
&= n\frac{K}{N}\\
&= 3\left(\frac{4}{9}\right)\\
&= \frac{4}{3}.
\end{align*}

The variance is

\begin{align*}
\operatorname{Var}(Y_1)
&= n\frac{K}{N}
\left(1-\frac{K}{N}\right)
\frac{N-n}{N-1}\\
&= 3\left(\frac{4}{9}\right)
\left(\frac{5}{9}\right)
\left(\frac{6}{8}\right)\\
&= \frac{5}{9}.
\end{align*}

\subsection{(b) Expectation and Variance of $Y_2$}

There are $K=3$ never-married executives among the
$N=9$ executives, and $n=3$ are selected.

The expected number of never-married executives is

\begin{align*}
E(Y_2)
&= n\frac{K}{N}\\
&= 3\left(\frac{3}{9}\right)\\
&= 1.
\end{align*}

The variance is

\begin{align*}
\operatorname{Var}(Y_2)
&= n\frac{K}{N}
\left(1-\frac{K}{N}\right)
\frac{N-n}{N-1}\\
&= 3\left(\frac{3}{9}\right)
\left(\frac{6}{9}\right)
\left(\frac{6}{8}\right)\\
&= \frac{1}{2}.
\end{align*}

\subsection{(c) Covariance Between $Y_1$ and $Y_2$}

By definition,

\[
\operatorname{Cov}(Y_1,Y_2)
= E(Y_1Y_2)-E(Y_1)E(Y_2).
\]

For two distinct categories in a multivariate
hypergeometric distribution,

\[
E(Y_1Y_2)
= \frac{n(n-1)K_1K_2}{N(N-1)},
\]

where $K_1=4$ is the number of married executives and
$K_2=3$ is the number of never-married executives.

Therefore,

\begin{align*}
E(Y_1Y_2)
&= \frac{3(3-1)(4)(3)}{9(9-1)}\\
&= \frac{72}{72}\\
&= 1.
\end{align*}

Using the expectations from parts (a) and (b),

\begin{align*}
\operatorname{Cov}(Y_1,Y_2)
&= E(Y_1Y_2)-E(Y_1)E(Y_2)\\
&= 1-\left(\frac{4}{3}\right)(1)\\
&= -\frac{1}{3}.
\end{align*}

The covariance is negative because selecting more married
executives leaves fewer positions available for
never-married executives in the sample.

\subsection{(d) Expectation and Variance of $Y_1+Y_2$}

Let $S=Y_1+Y_2$ denote the total number of married
and never-married executives selected.

By linearity of expectation,

\begin{align*}
E(S)
&= E(Y_1+Y_2)\\
&= E(Y_1)+E(Y_2)\\
&= \frac{4}{3}+1\\
&= \frac{7}{3}.
\end{align*}

For the variance of a sum of two random variables,

\[
\operatorname{Var}(Y_1+Y_2)
= \operatorname{Var}(Y_1)
+ \operatorname{Var}(Y_2)
+ 2\operatorname{Cov}(Y_1,Y_2).
\]

Substituting the results from parts (a), (b), and (c),

\begin{align*}
\operatorname{Var}(S)
&= \frac{5}{9}+\frac{1}{2}
+2\left(-\frac{1}{3}\right)\\
&= \frac{10}{18}+\frac{9}{18}
-\frac{12}{18}\\
&= \frac{7}{18}.
\end{align*}

As a verification, there are $4+3=7$ executives who
are either married or never-married. Thus,
$S\sim\operatorname{Hypergeometric}(9,7,3)$.

Applying the hypergeometric formulas directly,

\begin{align*}
E(S)
&= 3\left(\frac{7}{9}\right)
= \frac{7}{3},\\[6pt]
\operatorname{Var}(S)
&= 3\left(\frac{7}{9}\right)
\left(\frac{2}{9}\right)
\left(\frac{6}{8}\right)
= \frac{7}{18}.
\end{align*}

Both results agree with the calculations above.

\section{Problem 2 Answers}

We are given the joint probability density function

\[
f(y_1,y_2)=
\begin{cases}
4y_1y_2, & 0\leq y_1\leq 1,\ 0\leq y_2\leq 1,\\
0, & \text{elsewhere}.
\end{cases}
\]

First, we determine the marginal probability density
functions of $Y_1$ and $Y_2$.

For $Y_1$,

\begin{align*}
f_{Y_1}(y_1)
&= \int_0^1 f(y_1,y_2)\,dy_2\\
&= \int_0^1 4y_1y_2\,dy_2\\
&= 4y_1\left[\frac{y_2^2}{2}\right]_0^1\\
&= 2y_1,\qquad 0\leq y_1\leq 1.
\end{align*}

Similarly, for $Y_2$,

\begin{align*}
f_{Y_2}(y_2)
&= \int_0^1 f(y_1,y_2)\,dy_1\\
&= \int_0^1 4y_1y_2\,dy_1\\
&= 4y_2\left[\frac{y_1^2}{2}\right]_0^1\\
&= 2y_2,\qquad 0\leq y_2\leq 1.
\end{align*}

\subsection{(a) Expected Values of $Y_1$ and $Y_2$}

By definition, the expected value of a continuous
random variable $Y$ is

\[
E(Y)=\int_{-\infty}^{\infty} y f_Y(y)\,dy.
\]

Therefore,

\begin{align*}
E(Y_1)
&= \int_0^1 y_1 f_{Y_1}(y_1)\,dy_1\\
&= \int_0^1 y_1(2y_1)\,dy_1\\
&= 2\int_0^1 y_1^2\,dy_1\\
&= 2\left[\frac{y_1^3}{3}\right]_0^1\\
&= \frac{2}{3}.
\end{align*}

Similarly,

\begin{align*}
E(Y_2)
&= \int_0^1 y_2 f_{Y_2}(y_2)\,dy_2\\
&= \int_0^1 y_2(2y_2)\,dy_2\\
&= 2\int_0^1 y_2^2\,dy_2\\
&= 2\left[\frac{y_2^3}{3}\right]_0^1\\
&= \frac{2}{3}.
\end{align*}

Hence,

\[
E(Y_1)=\frac{2}{3},
\qquad
E(Y_2)=\frac{2}{3}.
\]

\subsection{(b) Variances of $Y_1$ and $Y_2$}

The variance of a random variable $Y$ is

\[
V(Y)=E(Y^2)-[E(Y)]^2.
\]

First, we calculate $E(Y_1^2)$:

\begin{align*}
E(Y_1^2)
&= \int_0^1 y_1^2 f_{Y_1}(y_1)\,dy_1\\
&= \int_0^1 y_1^2(2y_1)\,dy_1\\
&= 2\int_0^1 y_1^3\,dy_1\\
&= 2\left[\frac{y_1^4}{4}\right]_0^1\\
&= \frac{1}{2}.
\end{align*}

Therefore,

\begin{align*}
V(Y_1)
&= E(Y_1^2)-[E(Y_1)]^2\\
&= \frac{1}{2}-\left(\frac{2}{3}\right)^2\\
&= \frac{1}{2}-\frac{4}{9}\\
&= \frac{1}{18}.
\end{align*}

Similarly, for $Y_2$,

\begin{align*}
E(Y_2^2)
&= \int_0^1 y_2^2 f_{Y_2}(y_2)\,dy_2\\
&= 2\int_0^1 y_2^3\,dy_2\\
&= 2\left[\frac{y_2^4}{4}\right]_0^1\\
&= \frac{1}{2}.
\end{align*}

Thus,

\begin{align*}
V(Y_2)
&= E(Y_2^2)-[E(Y_2)]^2\\
&= \frac{1}{2}-\left(\frac{2}{3}\right)^2\\
&= \frac{1}{18}.
\end{align*}

Hence,

\[
V(Y_1)=\frac{1}{18},
\qquad
V(Y_2)=\frac{1}{18}.
\]

\subsection{(c) Expected Value of $Y_1-Y_2$}

Using the linearity of expectation,

\begin{align*}
E(Y_1-Y_2)
&= E(Y_1)-E(Y_2)\\
&= \frac{2}{3}-\frac{2}{3}\\
&= 0.
\end{align*}

Therefore,

\[
E(Y_1-Y_2)=0.
\]

\subsection{(d) Covariance Between $Y_1$ and $Y_2$}

By definition,

\[
\operatorname{Cov}(Y_1,Y_2)
= E(Y_1Y_2)-E(Y_1)E(Y_2).
\]

First, we calculate $E(Y_1Y_2)$ using the
joint probability density function:

\begin{align*}
E(Y_1Y_2)
&= \int_0^1\int_0^1
y_1y_2 f(y_1,y_2)\,dy_2\,dy_1\\
&= \int_0^1\int_0^1
4y_1^2y_2^2\,dy_2\,dy_1\\
&= 4\int_0^1 y_1^2
\left(\int_0^1 y_2^2\,dy_2\right)dy_1\\
&= 4\int_0^1 y_1^2
\left[\frac{y_2^3}{3}\right]_0^1dy_1\\
&= \frac{4}{3}\int_0^1 y_1^2\,dy_1\\
&= \frac{4}{3}
\left[\frac{y_1^3}{3}\right]_0^1\\
&= \frac{4}{9}.
\end{align*}

Therefore,

\begin{align*}
\operatorname{Cov}(Y_1,Y_2)
&= E(Y_1Y_2)-E(Y_1)E(Y_2)\\
&= \frac{4}{9}
-\left(\frac{2}{3}\right)\left(\frac{2}{3}\right)\\
&= \frac{4}{9}-\frac{4}{9}\\
&= 0.
\end{align*}

Hence,

\[
\operatorname{Cov}(Y_1,Y_2)=0.
\]

A covariance of zero indicates that $Y_1$ and $Y_2$
are uncorrelated, meaning that there is no linear
relationship between the two variables.

Although zero covariance does not necessarily imply
independence in general, we can establish independence
in this particular case.

Using the marginal density functions,

\begin{align*}
f_{Y_1}(y_1)f_{Y_2}(y_2)
&= (2y_1)(2y_2)\\
&= 4y_1y_2\\
&= f(y_1,y_2),
\end{align*}

for $0\leq y_1,y_2\leq 1$. Both sides are zero
outside this region.

Since the joint density is the product of the
marginal densities, $Y_1$ and $Y_2$ are independent.

\subsection{(e) Variance of $Y_1+Y_2$}

The variance of the sum of two random variables is

\[
V(Y_1+Y_2)
= V(Y_1)+V(Y_2)
+2\operatorname{Cov}(Y_1,Y_2).
\]

From parts (b) and (d), we have

\[
V(Y_1)=V(Y_2)=\frac{1}{18}
\]

and

\[
\operatorname{Cov}(Y_1,Y_2)=0.
\]

Therefore,

\begin{align*}
V(Y_1+Y_2)
&= V(Y_1)+V(Y_2)
+2\operatorname{Cov}(Y_1,Y_2)\\
&= \frac{1}{18}+\frac{1}{18}+2(0)\\
&= \frac{2}{18}\\
&= \frac{1}{9}.
\end{align*}

Hence,

\[
V(Y_1+Y_2)=\frac{1}{9}.
\]

\end{document}
